\section{An explicit spatial receiver for the original radial modes}
The positive residual-energy form admits a source-bound spatial incidence rather than an unspecified output weight. Retain the original regular quadrupole radial operator and normalized TT tensors. Its unitary spherical-Bessel transform is
\[
 (\mathcal H_2 f)(r)=\sqrt{\frac2\pi}\int_0^\infty k j_2(kr)f(k)\,\dd k.
\]
It maps $L^2(\dd k)$ unitarily to $L^2(r^2\dd r)$.
The identity $\int r^2j_2(kr)j_2(k'r)\dd r=\pi\delta(k-k')/(2k^2)$ fixes its normalization and retains the regular finite-energy extension at the origin. With the original angular isometry $\mathcal UQ=\sqrt{5/(8\pi)}\Lambda_nQ$, define
\[
 (Jf)(r,n)=\sum_{\alpha=+,\times}\mathcal Ue_\alpha(n)(\mathcal H_2 f_\alpha)(r).
\]
This $J$ is an isometry on the original two-polarization residual-energy output. A real bounded profile $0<N_-\le N(\mathbf x)\le N_+$ supplies $W_N=J^*NJ$ with the same positive bounds. The full closed form retains both cross terms and its entire counterterm.

This is a specified positive-frequency square-form spatial receiver. Identifying it with a pointwise canonical quantum stress tensor still requires the common local-action, vacuum, derivative-energy and boundary calculation. Its compressed domain does not silently add the full angular complement dynamics. For
\[
 A_{ab}(r)=\int N(r,n)\,\mathcal Ue_a:\mathcal Ue_b\,\dd\Omega,
\]
the spectral kernel is
\[
 W_{ab}(k,k')=\frac2\pi kk'\int_0^\infty
 r^2j_2(kr)j_2(k'r)A_{ab}(r)\dd r.
\]
A scalar spatial profile can therefore mix radial frequencies.

\section{Exact selection rules and the gradient complement}
The polarized TT contraction is
\[
 \begin{aligned}
 \Lambda Q:\Lambda Q'={}&\operatorname{tr}(QQ')
 -n^{\mathrm T}(QQ'+Q'Q)n\\
 &+\tfrac12(n^{\mathrm T}Qn)(n^{\mathrm T}Q'n).
 \end{aligned}
\]
Its spherical degrees are precisely among $L=0,2,4$. Thus every odd scalar harmonic and every scalar harmonic above degree four has zero compression. For the actual two driven tensors, write $C[h]_{ab}=\int h\,\mathcal Ue_a:\mathcal Ue_b\dd\Omega$. Exact integration gives
\[
\begin{gathered}
 C[1]=I,\quad C[x]=C[y]=C[z]=0,\\
 C[3z^2-1]=\tfrac47I,\quad C[P_4(z)]=\tfrac1{126}I,\\
 C[x^2-y^2]=C[2xy]=C[2xz]=C[2yz]=0,\\
 C[x^4-6x^2y^2+y^4]=\tfrac4{63}\begin{pmatrix}1&0\\0&-1\end{pmatrix},\\
 C[4xy(x^2-y^2)]=\tfrac4{63}\begin{pmatrix}0&1\\1&0\end{pmatrix}.
\end{gathered}
\]
The other degree-four components vanish. The real scalar angular image on the original pair is exactly the three-dimensional space of real symmetric matrices. On the entire five-STF image, the actual $L=0,2,4$ Gaunt map has rank fifteen. It realizes every small real symmetric five-coordinate weight with a bounded real scalar profile; positivity still constrains finite coefficients. The earlier imaginary antisymmetric polarization fixture belongs to a different realization class. Here $\{D,C\}=0$ and $D^2=C^2$ make the real weighted counterterm proportional to $D^2$, retaining all four original dark dimensions.

A vanishing compressed gradient does not remove its complement. With $P_5$ the full STF angular-image projector,
\[
 \begin{gathered}
 P_5zP_5=0,\\
 P_5z(I-P_5)zP_5\\
 =\operatorname{diag}(11/21,11/21,5/21,5/21,1/7).
 \end{gathered}
\]
The latter is positive and rank five, with rank two on the driven pair. A full angular evolution must retain this gradient leakage and its returning effect.

The global profile remains bounded and positive. Every polynomial Taylor expression below is restricted to the finite near ball; the actual bounded exterior remains in the far account. For a smooth source-centered profile, $N=N_0+g_ix_i+H_{ij}x_ix_j/2+\cdots$, its quadratic compression is
\[
 A(r)-N_0I=\frac{r^2}{42}(5\operatorname{tr}H+6H_{zz})I+\cdots.
\]
Cubic terms also cancel by parity. A fourth directional-derivative bound $M_4$ on $r\le L$ supplies
\[
 \delta_{\rm near}\le \frac{L^2}{42}|5\operatorname{tr}H+6H_{zz}|+\frac{M_4L^4}{24}.
\]
For $N=1-m/|\mathbf d+\mathbf x|+O(m^2/d^2)$, the leading term is $mr^2(1-3\cos^2\beta)/(7d^3)$. It vanishes at $\cos^2\beta=1/3$. Even a scalar polarization compression varies radially and need not become one common lapse for six different correlated preparations.

\section{An actual-state spatial tail and a finite profile criterion}
Let $K_{\rm in/out}[\Psi]=\|P_{\rm in/out}JT_g\Psi\|^2$ split the original complete positive form at radius $L$. The physical profile supplies
\[
 |r[\Phi,\Psi]|\le\delta_{\rm near}\sqrt{K_{\rm in}[\Phi]K_{\rm in}[\Psi]}
 +\delta_{\rm far}\sqrt{K_{\rm out}[\Phi]K_{\rm out}[\Psi]}.
\]
The exterior stays in this account. If $\|W-N_0I\|<N_0/2$, put $k_A=3K_{{\rm prep},\rm up}+4\varepsilon_{\rm up}$. A reference ceiling $K\le k_R$ and uniform tail ceiling $K_{\rm out}\le\eta_{\rm ref}(L)$ give
\[
 \|\Psi_{\rm act}(T)-\chi(T)\|^2
 \le2T\sqrt{k_A}\left[\delta_{\rm near}\sqrt{k_R}
                         +\delta_{\rm far}\sqrt{\eta_{\rm ref}(L)}\right].
\]
Four times the state-distance bound controls the actual current error. Comparing it to the retained nonzero reference amplitude supplies the same unwrapped arcsine phase criterion as before.

The tail can be bounded for the original exact preparation. The retained hard-branch pull-through identity gives, for $k\ge\lambda=64\Delta$,
\[
 a_\alpha(k)\Phi_j=g u_R(k)(H_h-E_j+k)^{-1}\sigma_\alpha\Phi_j,
 \quad H_h-E_j\ge-\Delta.
\]
For $H_j(k)=k(H_h-E_j+k)^{-1}$, spectral calculus gives $\|H_j\|\le64/63$ and $\|H_j'\|\le(1/4+64/3969)/k$. Applying the original Bessel transform to this Hilbert-space-valued function supplies an explicit conservative annihilation bound $\|A_\alpha(r)\|\le C_A(L)/r^2$ for $r\ge L\ge2R$. The retained derivation splits at $1/R$, integrates the low-frequency Bessel derivative envelope, and keeps all nine high-frequency sine/cosine products with phases $k(r\pm R)$.

The shift has the exact profile
\[
 F_\alpha(r)=\frac{g}{\pi\sqrt\kappa\,rR}
 Q_2\!\left(\frac{r^2+R^2}{2rR}\right)\sigma_\alpha\Phi,
\]
where its logarithmic shell singularity is square integrable. For $r\ge L$ it is bounded by $C_F(L)/r^4$, with
\[
 C_F(L)=\frac{16|g|R^2}{15\pi\sqrt\kappa(1-y_L^2)},\qquad
 y_L=\frac{2LR}{L^2+R^2}.
\]
Keeping both original branches, both polarizations and the common soft vacuum gives
\[
 K_{{\rm out},{\rm prep}}(L)\le
 \min\left\{K_{{\rm prep},\rm up},\frac{8C_A(L)^2}{L}
                              +\frac{4C_F(L)^2}{5L^5}\right\}.
\]
The original energy-norm Duhamel estimate $d_A^2\le8T^2g^2v_s(1+F_T+\lambda)$ then supplies
\[
 \eta_{\rm ref}(L)\le\min\{K_{{\rm prep},\rm up}+\varepsilon E(T),
               2K_{{\rm out},{\rm prep}}(L)+4(1+c_g)d_A^2\}.
\]
The comparison payment independent of $L$ is a bound feature, not an actual tail floor. At the illustrative radius $L=10^{15}R$ in model coordinates, the preparation-tail ceilings divided by the inherited $K_{\rm prep}$ ceilings are approximately $0.00275835$, $0.00166416$, $2.06463\times10^{-5}$, $0.00267798$, $2.23339\times10^{-5}$ and $2.43841\times10^{-5}$ in lane order 01,02,03,12,13,23. No SI ruler is assigned by these examples.

These receiving-form results supply an explicit incidence, realizable-weight classification, selection and complement certificates, and actual-state tail bounds. Appendix~\ref{app:tail-coefficient} displays the full standalone $C_A(L)$ coefficient. The Gaunt certificates do not evaluate the six time-dependent current-response integrals. The following six-lane continuation supplies one actual positive profile, a declared common dimensional ruler and evaluated response enclosures. It retains the isolated inhomogeneous two-time matrix elements as unresolved within those enclosures. These are kept separate from the local quantum stress identification.
